\chapter{Background}
\label{ch:background}

\section{Multicast DNS and DNS-Based Service Discovery}
\label{sec:mdns-dnssd}

I selected mDNS/DNS-SD because it implements the exact discovery patterns for networked peripherals I required for distributed AI services. Multicast DNS (mDNS), as defined in RFC 6762 \cite{rfc6762}, enables peer-to-peer name resolution on a local network by querying the multicast address \texttt{224.0.0.251} (or \texttt{ff02::fb} for IPv6) over port 5353. By utilizing the \texttt{.local} top-level domain, mDNS allows devices to resolve hostnames without reliance on external DNS infrastructure, fulfilling the "zero-configuration" vision described by Guttman \cite{guttman2001zeroconf}. DNS-based Service Discovery (DNS-SD), specified in RFC 6763 \cite{rfc6763}, layers structured advertisement on top of mDNS through three DNS record types: PTR records enumerate service instances, SRV records provide the hostname and port, and TXT records carry key-value metadata. 

The viability of Saturn relies on the broad cross-platform support of these protocols. Apple's Bonjour \cite{bonjour}, released in 2002, ships natively with macOS and iOS and is available as a Windows service, while Avahi \cite{avahi} serves as the default mDNS daemon on most Linux distributions. As Siddiqui et al. \cite{siddiqui2012zeroconf} confirmed, Avahi and Bonjour interoperate correctly on typical local networks, and Siljanovski et al. \cite{siljanovski2014mdns} demonstrated that mDNS functions even on resource-constrained hardware. Together, these implementations allow Saturn to leverage OS-level discovery without requiring users to install external daemons, maintaining the project's zero-configuration premise.

By leveraging TXT records, Saturn encapsulates structured metadata—including endpoint URLs, API types (such as OpenAI-compatible or Ollama), and authentication credentials—into a single announcement. This effectively automates the manual configuration typically required for AI services. However, this approach necessitates designing around specific mDNS constraints: a 255-byte limit per TXT string, link-local scope restrictions, and the inherent lack of native broadcast security. 

The same broadcast visibility that enables discovery also creates a privacy surface. Kaiser and Waldvogel \cite{kaiser2014privacy} showed that passive eavesdropping on mDNS traffic reveals which services a device offers and consumes, while Könings et al. \cite{konings2013devicenames} found that 59\% of mDNS device names contain the owner's real name. For Saturn, these findings mean that AI endpoint announcements—including ephemeral API keys—are visible to every device on the local network. To mitigate these risks, Chapter \ref{ch:design} details the specific TXT schema and the use of ephemeral keys, while Chapter \ref{ch:discussion} provides a full evaluation of the protocol's performance and security trade-offs at scale.protocol's performance at scale.

\section{Alternative Protocol Analysis}
\label{sec:comparisons}

I evaluated four alternative service discovery protocols against three core requirements: cross-platform support, structured metadata, and zero-infrastructure operation. Among these, NetBIOS \cite{rfc1001} is the least viable, as it lacks structured metadata beyond a 15-character hostname and has been deprecated by Microsoft in favor of mDNS since Windows 10. Similarly, DLNA \cite{dlna2017} is unsuitable because it restricts its scope to media streaming, its consortium dissolved in 2017, and its schema cannot express critical information like API types or authentication tokens.

WS-Discovery \cite{wsdiscovery} and UPnP \cite{upnp2015} offer more robust metadata capabilities but introduce significant overhead. WS-Discovery relies on SOAP/XML messages, where a single probe consumes roughly 1\,KB of XML compared to a 50-byte DNS query; such verbosity makes Saturn beacons impractical for weaker networks. UPnP comes closest to meeting the requirements due to its rich device descriptions and active deployment, yet it bundles unnecessary control features like event subscriptions. Furthermore, security vulnerabilities in UPnP's control surface have led US-CERT to recommend its deactivation, whereas Saturn requires lean discovery without remote-control risks.

The structural insight of the Dynamic Host Configuration Protocol (DHCP) \cite{rfc2131} served as the primary inspiration for Saturn's design. In a DHCP workflow, a device broadcasts a DISCOVER message and receives an OFFER, allowing it to operate without manual user configuration. Saturn adopts this concentration of complexity to eliminate per-user configuration for AI service endpoints. By centralizing the logic within the protocol, one administrator can configure the server while every client benefits automatically, validating the principle of infrastructure-level provisioning. Despite this conceptual alignment, Saturn departs from DHCP in three functional ways that necessitated a new implementation. First, Saturn manages multiple dynamic services as they start and stop, rather than assigning a single static address at network join. Second, Saturn operates entirely in userspace without requiring the privileged access necessary to configure network interfaces. Finally, Saturn communicates via multicast DNS on port 5353 after connectivity is established, rather than operating below the IP layer on UDP port 67. These distinctions ensure that Saturn provides high-level AI service discovery while maintaining the zero-configuration strengths of the DHCP model.